\section*{Statement on the Use of Generative AI}

\begingroup
\small
\setlength{\parskip}{0.35em}

During the preparation of this thesis, generative artificial intelligence was
used at different stages for a limited set of supporting tasks. The tool used
was OpenAI Codex with the GPT-5.6 Sol large language model.

Generative AI was used for the following purposes:

\begin{itemize}[leftmargin=1.5em,itemsep=0.1em,topsep=0.25em,parsep=0pt]
  \item correcting spelling, grammar, and orthographic errors in existing text;
  \item improving the readability and flow of existing text without
        introducing new substantive arguments;
  \item assisting in the discovery of potentially relevant research
        publications;
  \item supporting the formatting of text, headings, formulas, and other
        document elements based on specifications or prototypes provided by
        the author;
  \item helping to ensure consistent use of mathematical notation and symbols
        throughout the thesis;
  \item generating and refining LaTeX-native graphs and illustrations from
        data, requirements, and visual prototypes provided by the author;
  \item performing one limited data-normalization task for the heatmap
        presented in Appendix~B, in which related category labels were
        consolidated to reduce the number of rows and columns; and
  \item assisting with occasional technical debugging and with the
        implementation of specified procedures in R and in the FELICS
        reference implementation.
\end{itemize}

All methodological decisions, including the conceptualization and development
of the FELICS framework, the experimental design, and the execution of the
experiments, were made and carried out by the author. Generative AI was not used
to:

\begin{itemize}[leftmargin=1.5em,itemsep=0.1em,topsep=0.25em,parsep=0pt]
  \item generate original substantive academic prose;
  \item serve as a source of scientific evidence or extract findings from
        research publications beyond providing an initial overview;
  \item conceptualize or develop the FELICS framework;
  \item make methodological decisions or independently design or conduct the
        experiments; or
  \item formulate interpretations, conclusions, or substantive scientific
        claims.
\end{itemize}

All AI-assisted outputs were independently reviewed and verified by the author.
Generated code was examined manually and tested through execution. Publications
suggested through AI assistance were checked against the original sources.
Numerical and graphical outputs were compared with their underlying
experimental data. Where appropriate, independently generated visualizations---
for example, plots produced in R directly from the experimental data---were used
to verify the values and structure of AI-assisted LaTeX-native plots.

No confidential, personal, or otherwise sensitive data, and no non-public
third-party copyrighted material, was submitted to the external AI service.
Prompts and individual instances of AI use are not reproduced because
generative AI was employed repeatedly for the narrowly defined supporting tasks
listed above rather than as a source of academic content.

The author assumes full responsibility for the accuracy, integrity, and final
content of this thesis, including its text, sources, software, analyses,
figures, interpretations, and conclusions.

\endgroup
